\begin{tikzpicture}[
    line cap=round,
    line join=round,
    >=Latex,
    font=\footnotesize,
    factor/.style={
        circle,
        draw=black!68,
        fill=white,
        minimum size=8.0mm,
        inner sep=0pt,
        line width=0.65pt,
        text=black!82
    },
    core/.style={
        circle,
        draw=black!68,
        fill=black!7,
        minimum size=8.0mm,
        inner sep=0pt,
        line width=0.65pt,
        text=black!82
    },
    edge/.style={
        draw=black!68,
        line width=0.65pt
    },
    hedge/.style={
        draw=black!43,
        dashed,
        line width=0.55pt,
        dash pattern=on 3.0pt off 2.4pt
    },
    arrow/.style={
        text=black!45,
        font=\large
    },
    captionlabel/.style={
        font=\footnotesize,
        text=black!78
    }
]

% Same visual grammar as the original tensor-network sketch, but with lighter
% strokes and labels aligned with the explicit p-fold Tucker notation.

% =========================================================
% Tucker  (3D-like view)
% =========================================================
\begin{scope}[xshift=0cm]
    \node[factor] (f1) at (0.00, 0.00) {$\mathbf{A}^{(1)}$};
    \node[factor] (f2) at (3.464, 2.000) {$\mathbf{A}^{(2)}$};
    \node[factor] (f3) at (2.249, 2.932) {$\mathbf{A}^{(3)}$};

    \node[core] (c1) at (1.732, 1.00) {$\mathcal{X}$};

    \draw[edge] (f1) -- (c1);
    \draw[edge] (f2) -- (c1);
    \draw[edge] (f3) -- (c1);

    \node[captionlabel] at (1.732,-1.05) {Tucker \((p=1)\)};
\end{scope}

\node[arrow] at (4.00, 1.00) {$\Longrightarrow$};

% =========================================================
% Quadratic Tucker
% =========================================================
\begin{scope}[xshift=5cm]
    \node[factor] (f1) at (0.00, 0.00) {$\mathcal{A}^{(1)}$};
    \node[factor] (f2) at (3.464, 2.000) {$\mathcal{A}^{(2)}$};
    \node[factor] (f3) at (2.249, 2.932) {$\mathcal{A}^{(3)}$};

    \node[core] (c1) at (1.378, 1.354) {$\mathcal{X}$};
    \node[core] (c2) at (2.085, 0.64) {$\mathcal{X}$};

    \draw[edge] (f1) -- (c1);
    \draw[edge] (f1) -- (c2);
    \draw[hedge] (f2) -- (c1);
    \draw[edge] (f2) -- (c2);
    \draw[edge] (f3) -- (c1);
    \draw[edge] (f3) -- (c2);

    \node[captionlabel] at (1.732,-1.05) {Quadratic \((p=2)\)};
\end{scope}

\node[arrow] at (9.00, 1.00) {$\Longrightarrow$};

% =========================================================
% p-Fold Tucker
% =========================================================
\begin{scope}[xshift=10cm]
    \node[factor] (f1) at (0.00, 0.00) {$\mathcal{A}^{(1)}$};
    \node[factor] (f2) at (3.464, 2.000) {$\mathcal{A}^{(2)}$};
    \node[factor] (f3) at (2.249, 2.932) {$\mathcal{A}^{(3)}$};

    \node[core] (c1) at (1.378, 1.354) {$\mathcal{X}$};
    \node[core] (c3) at (0.671, 2.061) {$\mathcal{X}$};
    \node[core] (c4) at (2.793, -0.061) {$\mathcal{X}$};

    \draw[edge] (f1) -- (c1);
    \draw[edge] (f1) -- (c3);
    \draw[edge] (f1) -- (c4);
    \draw[hedge] (f2) -- (c1);
    \draw[hedge] (f2) -- (c3);
    \draw[edge] (f2) -- (c4);
    \draw[edge] (f3) -- (c1);
    \draw[edge] (f3) -- (c3);
    \draw[edge] (f3) -- (c4);

    \node[font=\large, rotate=-45, text=black!55] at (2.085, 0.64) {$\cdots$};

    \node[captionlabel] at (1.732,-1.05) {Formal \(p\)-fold};
\end{scope}

\end{tikzpicture}
